\section{Introduction}

This report documents a complete small-language-model program carried out
on TPU Research Cloud hardware between 2026-07-20 and 2026-08-08: kernel
development on a single-host TPU v6e-8, a 367B-token pretraining campaign
on a 32-chip TPU v4-64, and a post-training stage (supervised fine-tuning
plus offline logit distillation) evaluated against the strongest publicly
available models of the same size class, all measured under one harness.

The program's through-line is that every layer of the stack was
co-designed with the ones above and below it. The architecture's gate
bound is also the kernel's floating-point overflow budget
(\S\ref{sec:kernels}). The tokenizer is not merely small; it is an exact
subset projection of a 4B teacher's vocabulary, which later deletes an
entire class of distillation alignment machinery
(\S\ref{sec:tokenizer}, \S\ref{sec:posttraining}). The optimizer's
logit-growth guard is priced at zero because its trigger never fired in
700{,}000 steps --- a fact we can state only because the training loop
recorded per-head maximum logits the whole way (\S\ref{sec:optimizer}).

\paragraph{Contributions.}
\begin{enumerate}
  \item A hybrid gated-delta-attention (KDA) model,
  \texttt{kda\_hybrid\_yx49k\_l20}: 307{,}698{,}680 parameters, 20 layers
  as five \([\mathrm{KDA},\mathrm{KDA},\mathrm{KDA},\mathrm{GQA}]\)
  cycles, trained at sequence length 2048 (\S\ref{sec:architecture}).
  \item Fused TPU training kernels for KDA built on a chunkwise WY form:
  a divide-and-conquer explicit inverse that is exact by nilpotency, a
  decay-anchored dot-product factorization with a provable fp32 overflow
  bound, and a split-backward variant for the older v4 ISA
  (\S\ref{sec:kernels}), with a numerics policy qualified against fp64 on
  real text (\S\ref{sec:numerics}).
  \item A matrix/scalar split optimizer --- orthogonalized momentum for
  matrices, AdamW for the rest --- with a grouped-query adaptation of
  attention-logit clipping (\S\ref{sec:optimizer}).
  \item The yx49k tokenizer: 49{,}152 tokens, each mapping to a distinct
  token of the teacher vocabulary, at practical parity with a
  2.6$\times$ larger reference tokenizer on bytes-per-token
  (\S\ref{sec:tokenizer}).
  \item The campaign itself: 700k steps at 524{,}288 tokens per step on
  32 chips, median 1.09M tokens/s, a mid-anneal crash with a
  measurably clean resume, and a terminal anneal worth $-0.195$ nats of
  holdout loss (\S\ref{sec:pretraining}, \S\ref{sec:evaluation}).
  \item A post-training arc --- SFT at 6.27M rows, then offline top-$K$
  logit distillation from a frozen 4B teacher --- with every design
  decision A/B-measured, including a clean null result on repeated
  epochs (\S\ref{sec:posttraining}), and an external calibration against
  Gemma~3 270M IT and SmolLM 360M v1/v2 under identical measurement
  (\S\ref{sec:calibration}).
\end{enumerate}

\begin{table}[h]
\centering
\small
\caption{Summary card.}
\label{tab:summary}
\begin{tabular}{ll}
\toprule
Parameters & 307{,}698{,}680 (tied embeddings) \\
Layers & 20 = 5 $\times$ [KDA, KDA, KDA, GQA], $d_{\mathrm{model}}=1024$ \\
Tokenizer & yx49k (49{,}152 BPE, exact teacher subset) \\
Pretraining data & ClimbMix, 367B step-tokens ($\sim$341B unique) \\
Hardware & TPU v4-64: 32 chips / 8 hosts, pure data parallel \\
Throughput & median 1.09M tokens/s (480 ms / 524{,}288-token step) \\
Final loss & train 2.3617, holdout 2.3494 (ppl 10.48) \\
0-shot & HellaSwag 56.9, PIQA 75.2, ARC-e 67.5, WinoGrande 59.4 \\
Post-training & IFEval mean 51.0 (GOLD mix300k checkpoint) \\
\bottomrule
\end{tabular}
\end{table}
